% !TEX root = ../math.tex
\section{Byte-Pair Encoding Mechanics}
\label{sec:bpe-mechanics}
\label{sec:appendix-bpe-mechanics}

We now specialize the general tokenization formalism of this book to byte-pair encoders.
Our goal is to prove a structural property that is stronger than boundary sufficiency and is
specific to BPE tokenizers.

Throughout this section we identify each BPE token with its spelling.
Thus a token is itself a nonempty string over the character alphabet $\Lambda$,
and the literalization map is just concatenation of tokens.

\begin{definition}[Character Tokenization]
    For a string $x = x_1 \cdots x_n$, let
    \[
    \operatorname{chars}(x) = [x_1, \ldots, x_n]
    \]
    denote the sequence of one-character strings whose concatenation is $x$.
\end{definition}

\begin{definition}[BPE Merge System]
    A BPE merge system of depth $N$ consists of vocabularies
    \[
    \boldsymbol{\Lambda}^{(0)} \subseteq \boldsymbol{\Lambda}^{(1)} \subseteq \cdots \subseteq \boldsymbol{\Lambda}^{(N)}
    \]
    and tokens $\boldsymbol{\alpha_k}, \boldsymbol{\beta_k}, \boldsymbol{\gamma_k}$ for each
    $1 \leq k \leq N$ satisfying:
    \begin{enumerate}
        \item $\boldsymbol{\Lambda}^{(0)} = \Lambda$, where each character is regarded as a one-character token.
        \item $\boldsymbol{\alpha_k}, \boldsymbol{\beta_k} \in \boldsymbol{\Lambda}^{(k-1)}$.
        \item $\boldsymbol{\gamma_k} = \boldsymbol{\alpha_k} + \boldsymbol{\beta_k}$ as strings.
        \item $\boldsymbol{\gamma_k} \notin \boldsymbol{\Lambda}^{(k-1)}$.
        \item $\boldsymbol{\Lambda}^{(k)} = \boldsymbol{\Lambda}^{(k-1)} \cup \{\boldsymbol{\gamma_k}\}$.
    \end{enumerate}
\end{definition}

\begin{definition}[Single BPE Merge Pass]
    For $1 \leq k \leq N$, the $k$th merge pass
    \[
    M_k : (\boldsymbol{\Lambda}^{(k-1)})^* \to (\boldsymbol{\Lambda}^{(k)})^*
    \]
    is defined recursively by
    \[
    M_k([]) = [],
    \]
    \[
    M_k([\boldsymbol{\alpha_k}, \boldsymbol{\beta_k}] + \mathbf{r})
    =
    [\boldsymbol{\gamma_k}] + M_k(\mathbf{r}),
    \]
    and
    \[
    M_k([\boldsymbol{\tau}] + \mathbf{r})
    =
    [\boldsymbol{\tau}] + M_k(\mathbf{r})
    \]
    whenever either $\mathbf{r} = []$, or the first token of $\mathbf{r}$ is not $\boldsymbol{\beta_k}$,
    or $\boldsymbol{\tau} \neq \boldsymbol{\alpha_k}$.

    Thus $M_k$ performs the left-to-right maximal non-overlapping replacement of
    $[\boldsymbol{\alpha_k}, \boldsymbol{\beta_k}]$ by $[\boldsymbol{\gamma_k}]$.
\end{definition}

\begin{definition}[Stagewise BPE Tokenizers]
    The tokenizer after $k$ merge stages is the map
    \[
    T_k : S \to (\boldsymbol{\Lambda}^{(k)})^*
    \]
    defined recursively by
    \[
    T_0(x) = \operatorname{chars}(x),
    \qquad
    T_k(x) = M_k(T_{k-1}(x)).
    \]
    The final BPE tokenizer is $T := T_N$.
\end{definition}

\begin{definition}[Stagewise Canonicality]
    A token sequence $\mathbf{a} \in (\boldsymbol{\Lambda}^{(k)})^*$ is stage-$k$ canonical if and only if
    \[
    T_k(T^{-1}(\mathbf{a})) = \mathbf{a}.
    \]
    We denote the set of stage-$k$ canonical token sequences by $\mathbf{C}^{(k)}$.
\end{definition}

\begin{proposition}[Stage-$0$ Canonicality]
    \label{prop:bpe-stage-zero-canonicality}
    A token sequence is stage-$0$ canonical if and only if each of its tokens has length $1$.
\end{proposition}
\begin{proof}
    By definition,
    \[
    T_0(T^{-1}(\mathbf{a})) = \operatorname{chars}(T^{-1}(\mathbf{a})).
    \]
    The token sequence on the right is obtained by splitting the literalization of $\mathbf{a}$
    into one-character strings.
    Therefore it is equal to $\mathbf{a}$ if and only if every token of $\mathbf{a}$ is already
    a one-character string.
\end{proof}

\begin{definition}[Single-Stage Decompression]
    For $1 \leq k \leq N$, the single-stage decompression map
    \[
    D_k : (\boldsymbol{\Lambda}^{(k)})^* \to (\boldsymbol{\Lambda}^{(k-1)})^*
    \]
    is defined recursively by
    \[
    D_k([]) = [],
    \]
    \[
    D_k([\boldsymbol{\gamma_k}] + \mathbf{r})
    =
    [\boldsymbol{\alpha_k}, \boldsymbol{\beta_k}] + D_k(\mathbf{r}),
    \]
    and
    \[
    D_k([\boldsymbol{\tau}] + \mathbf{r})
    =
    [\boldsymbol{\tau}] + D_k(\mathbf{r})
    \qquad
    \text{for } \boldsymbol{\tau} \neq \boldsymbol{\gamma_k}.
    \]
\end{definition}

\begin{definition}[Stage-$k$ Reduced Sequence]
    For $1 \leq k \leq N$, we say that a token sequence $\mathbf{a}$ is stage-$k$ reduced if and only if
    it contains no adjacent occurrence of the pair $[\boldsymbol{\alpha_k}, \boldsymbol{\beta_k}]$.
    Equivalently,
    \[
    \operatorname{red}_k(\mathbf{a})
    \iff
    \not\exists i < |\mathbf{a}|-1,\;
    \mathbf{a}[i] = \boldsymbol{\alpha_k}
    \land
    \mathbf{a}[i+1] = \boldsymbol{\beta_k}.
    \]
\end{definition}

\begin{lemma}[Single-stage decompression preserves literalization]
    \label{lemma:bpe-decompression-preserves-literalization}
    For every $\mathbf{a} \in (\boldsymbol{\Lambda}^{(k)})^*$,
    \[
    T^{-1}(D_k(\mathbf{a})) = T^{-1}(\mathbf{a}).
    \]
\end{lemma}
\begin{proof}
    We argue by induction on $|\mathbf{a}|$.

    If $\mathbf{a} = []$, the claim is immediate.
    Otherwise write $\mathbf{a} = [\boldsymbol{\tau}] + \mathbf{r}$.

    If $\boldsymbol{\tau} = \boldsymbol{\gamma_k}$, then
    \[
    D_k(\mathbf{a}) = [\boldsymbol{\alpha_k}, \boldsymbol{\beta_k}] + D_k(\mathbf{r}),
    \]
    so
    \begin{align*}
    T^{-1}(D_k(\mathbf{a}))
    &=
    \boldsymbol{\alpha_k} + \boldsymbol{\beta_k} + T^{-1}(D_k(\mathbf{r})) \\
    &=
    \boldsymbol{\gamma_k} + T^{-1}(\mathbf{r}) \\
    &=
    T^{-1}(\mathbf{a}),
    \end{align*}
    using the defining relation $\boldsymbol{\gamma_k} = \boldsymbol{\alpha_k}+\boldsymbol{\beta_k}$
    and the inductive hypothesis.

    If $\boldsymbol{\tau} \neq \boldsymbol{\gamma_k}$, then
    \[
    D_k(\mathbf{a}) = [\boldsymbol{\tau}] + D_k(\mathbf{r}),
    \]
    and therefore
    \[
    T^{-1}(D_k(\mathbf{a}))
    =
    \boldsymbol{\tau} + T^{-1}(D_k(\mathbf{r}))
    =
    \boldsymbol{\tau} + T^{-1}(\mathbf{r})
    =
    T^{-1}(\mathbf{a}),
    \]
    again by the inductive hypothesis.
\end{proof}

\begin{lemma}[Single-stage decompression respects concatenation]
    \label{lemma:bpe-decompression-append}
    For all token sequences $\mathbf{a}$ and $\mathbf{b}$,
    \[
    D_k(\mathbf{a}+\mathbf{b}) = D_k(\mathbf{a}) + D_k(\mathbf{b}).
    \]
\end{lemma}
\begin{proof}
    We argue by induction on $|\mathbf{a}|$.

    If $\mathbf{a} = []$, the claim is immediate.
    Otherwise write $\mathbf{a} = [\boldsymbol{\tau}] + \mathbf{r}$ and split on whether
    $\boldsymbol{\tau} = \boldsymbol{\gamma_k}$.
    In each case the recursive definition of $D_k$ reduces the claim to the inductive hypothesis.
\end{proof}

\begin{lemma}[Decompression inverts a merge pass]
    \label{lemma:bpe-decompression-left-inverse}
    For every $\mathbf{a} \in (\boldsymbol{\Lambda}^{(k-1)})^*$,
    \[
    D_k(M_k(\mathbf{a})) = \mathbf{a}.
    \]
\end{lemma}
\begin{proof}
    We argue by induction on $|\mathbf{a}|$.

    If $\mathbf{a} = []$, the claim is immediate.
    Otherwise write $\mathbf{a} = [\boldsymbol{\tau}] + \mathbf{r}$.

    If $\boldsymbol{\tau} = \boldsymbol{\alpha_k}$ and $\mathbf{r} = [\boldsymbol{\beta_k}] + \mathbf{r}'$,
    then
    \[
    M_k(\mathbf{a}) = [\boldsymbol{\gamma_k}] + M_k(\mathbf{r}'),
    \]
    so
    \[
    D_k(M_k(\mathbf{a}))
    =
    [\boldsymbol{\alpha_k}, \boldsymbol{\beta_k}] + D_k(M_k(\mathbf{r}'))
    =
    [\boldsymbol{\alpha_k}, \boldsymbol{\beta_k}] + \mathbf{r}'
    =
    \mathbf{a},
    \]
    by the inductive hypothesis.

    Otherwise,
    \[
    M_k(\mathbf{a}) = [\boldsymbol{\tau}] + M_k(\mathbf{r}),
    \]
    and since $\boldsymbol{\tau} \neq \boldsymbol{\gamma_k}$,
    \[
    D_k(M_k(\mathbf{a}))
    =
    [\boldsymbol{\tau}] + D_k(M_k(\mathbf{r}))
    =
    [\boldsymbol{\tau}] + \mathbf{r}
    =
    \mathbf{a},
    \]
    again by the inductive hypothesis.
\end{proof}

\begin{lemma}[Merge pass yields a reduced sequence]
    \label{lemma:bpe-merge-produces-reduced}
    For every $\mathbf{a} \in (\boldsymbol{\Lambda}^{(k-1)})^*$,
    \[
    \operatorname{red}_k(M_k(\mathbf{a})).
    \]
\end{lemma}
\begin{proof}
    We argue by induction on $|\mathbf{a}|$.

    If $\mathbf{a} = []$, the claim is immediate.
    Otherwise write $\mathbf{a} = [\boldsymbol{\tau}] + \mathbf{r}$.

    If $\boldsymbol{\tau} = \boldsymbol{\alpha_k}$ and $\mathbf{r} = [\boldsymbol{\beta_k}] + \mathbf{r}'$,
    then
    \[
    M_k(\mathbf{a}) = [\boldsymbol{\gamma_k}] + M_k(\mathbf{r}').
    \]
    By the inductive hypothesis, the suffix $M_k(\mathbf{r}')$ is stage-$k$ reduced.
    Since the first token is $\boldsymbol{\gamma_k}$ rather than $\boldsymbol{\alpha_k}$,
    the whole sequence is stage-$k$ reduced.

    Otherwise,
    \[
    M_k(\mathbf{a}) = [\boldsymbol{\tau}] + M_k(\mathbf{r}).
    \]
    By the inductive hypothesis, the suffix $M_k(\mathbf{r})$ is stage-$k$ reduced.
    It remains only to rule out an adjacent occurrence of
    $[\boldsymbol{\alpha_k}, \boldsymbol{\beta_k}]$ starting at the first token.

    If $\boldsymbol{\tau} \neq \boldsymbol{\alpha_k}$, this is immediate.
    If $\boldsymbol{\tau} = \boldsymbol{\alpha_k}$, then by hypothesis the first token of $\mathbf{r}$
    is either absent or different from $\boldsymbol{\beta_k}$.
    Hence the first token of $M_k(\mathbf{r})$ is also either absent or different from $\boldsymbol{\beta_k}$.
    Therefore no adjacent occurrence of $[\boldsymbol{\alpha_k}, \boldsymbol{\beta_k}]$
    starts at the first token.
\end{proof}

\begin{lemma}[Reduced sequences are reconstructed by decompression]
    \label{lemma:bpe-reduced-reconstruction}
    If $\operatorname{red}_k(\mathbf{a})$, then
    \[
    M_k(D_k(\mathbf{a})) = \mathbf{a}.
    \]
\end{lemma}
\begin{proof}
    We argue by induction on $|\mathbf{a}|$.

    If $\mathbf{a} = []$, the claim is immediate.
    Otherwise write $\mathbf{a} = [\boldsymbol{\tau}] + \mathbf{r}$.

    If $\boldsymbol{\tau} = \boldsymbol{\gamma_k}$, then
    \[
    D_k(\mathbf{a}) = [\boldsymbol{\alpha_k}, \boldsymbol{\beta_k}] + D_k(\mathbf{r}),
    \]
    and hence
    \[
    M_k(D_k(\mathbf{a}))
    =
    [\boldsymbol{\gamma_k}] + M_k(D_k(\mathbf{r}))
    =
    [\boldsymbol{\gamma_k}] + \mathbf{r}
    =
    \mathbf{a},
    \]
    by the inductive hypothesis.

    If $\boldsymbol{\tau} \neq \boldsymbol{\gamma_k}$, then
    \[
    D_k(\mathbf{a}) = [\boldsymbol{\tau}] + D_k(\mathbf{r}).
    \]
    Because $\operatorname{red}_k(\mathbf{a})$, either $\boldsymbol{\tau} \neq \boldsymbol{\alpha_k}$,
    or $\mathbf{r} = []$, or the first token of $\mathbf{r}$ is not $\boldsymbol{\beta_k}$.
    In all cases the first clause in the definition of $M_k$ does not apply to
    $[\boldsymbol{\tau}] + D_k(\mathbf{r})$, so
    \[
    M_k(D_k(\mathbf{a}))
    =
    [\boldsymbol{\tau}] + M_k(D_k(\mathbf{r}))
    =
    [\boldsymbol{\tau}] + \mathbf{r}
    =
    \mathbf{a},
    \]
    again by the inductive hypothesis.
\end{proof}

\begin{theorem}[Stagewise canonicality characterization]
    \label{thm:bpe-canonical-characterization}
    Let $1 \leq k \leq N$.
    Then for every token sequence $\mathbf{a} \in (\boldsymbol{\Lambda}^{(k)})^*$,
    \[
    \mathbf{a} \in \mathbf{C}^{(k)}
    \iff
    \operatorname{red}_k(\mathbf{a})
    \land
    D_k(\mathbf{a}) \in \mathbf{C}^{(k-1)}.
    \]
\end{theorem}
\begin{proof}
    We prove both implications.

    First suppose $\mathbf{a} \in \mathbf{C}^{(k)}$.
    Then
    \[
    \mathbf{a}
    =
    T_k(T^{-1}(\mathbf{a}))
    =
    M_k(T_{k-1}(T^{-1}(\mathbf{a}))).
    \]
    Therefore $\operatorname{red}_k(\mathbf{a})$ by lemma~\ref{lemma:bpe-merge-produces-reduced}.

    Also,
    \[
    D_k(\mathbf{a})
    =
    D_k(M_k(T_{k-1}(T^{-1}(\mathbf{a}))))
    =
    T_{k-1}(T^{-1}(\mathbf{a}))
    \]
    by lemma~\ref{lemma:bpe-decompression-left-inverse}. By
    lemma~\ref{lemma:bpe-decompression-preserves-literalization},
    \[
    T^{-1}(D_k(\mathbf{a})) = T^{-1}(\mathbf{a}),
    \]
    so
    \[
    T_{k-1}(T^{-1}(D_k(\mathbf{a}))) = D_k(\mathbf{a}).
    \]
    Hence $D_k(\mathbf{a}) \in \mathbf{C}^{(k-1)}$.

    Conversely, suppose
    \[
    \operatorname{red}_k(\mathbf{a})
    \qquad\text{and}\qquad
    D_k(\mathbf{a}) \in \mathbf{C}^{(k-1)}.
    \]
    Then
    \[
    T_{k-1}(T^{-1}(D_k(\mathbf{a}))) = D_k(\mathbf{a}).
    \]
    By lemma~\ref{lemma:bpe-decompression-preserves-literalization},
    \[
    T^{-1}(D_k(\mathbf{a})) = T^{-1}(\mathbf{a}),
    \]
    so
    \[
    T_k(T^{-1}(\mathbf{a}))
    =
    M_k(T_{k-1}(T^{-1}(\mathbf{a})))
    =
    M_k(D_k(\mathbf{a})).
    \]
    Since $\operatorname{red}_k(\mathbf{a})$, lemma~\ref{lemma:bpe-reduced-reconstruction} gives
    \[
    M_k(D_k(\mathbf{a})) = \mathbf{a}.
    \]
    Therefore
    \[
    T_k(T^{-1}(\mathbf{a})) = \mathbf{a},
    \]
    i.e. $\mathbf{a} \in \mathbf{C}^{(k)}$.
\end{proof}

\begin{theorem}[BPE overlap concatenation theorem]
    \label{thm:bpe-overlap-concatenation}
    For every $k \in \{0,\ldots,N\}$, if
    \[
    \mathbf{b} \neq \boldsymbol{\epsilon},
    \qquad
    \mathbf{a}+\mathbf{b} \in \mathbf{C}^{(k)}
    \qquad\text{and}\qquad
    \mathbf{b}+\mathbf{c} \in \mathbf{C}^{(k)},
    \]
    then
    \[
    \mathbf{a}+\mathbf{b}+\mathbf{c} \in \mathbf{C}^{(k)}.
    \]
\end{theorem}
\begin{proof}
    We argue by induction on $k$.

    For $k=0$, proposition~\ref{prop:bpe-stage-zero-canonicality} implies that every token
    occurring in $\mathbf{a}+\mathbf{b}$ and every token occurring in $\mathbf{b}+\mathbf{c}$
    has length $1$. Hence every token in $\mathbf{a}+\mathbf{b}+\mathbf{c}$ has length $1$,
    so $\mathbf{a}+\mathbf{b}+\mathbf{c} \in \mathbf{C}^{(0)}$ by the same proposition.

    Now let $1 \leq k \leq N$, and suppose the theorem has been proved for $k-1$.
    Assume
    \[
    \mathbf{b} \neq \boldsymbol{\epsilon},
    \qquad
    \mathbf{a}+\mathbf{b} \in \mathbf{C}^{(k)}
    \qquad\text{and}\qquad
    \mathbf{b}+\mathbf{c} \in \mathbf{C}^{(k)}.
    \]
    By theorem~\ref{thm:bpe-canonical-characterization},
    \[
    \operatorname{red}_k(\mathbf{a}+\mathbf{b}),
    \quad
    D_k(\mathbf{a}+\mathbf{b}) \in \mathbf{C}^{(k-1)},
    \quad
    \operatorname{red}_k(\mathbf{b}+\mathbf{c}),
    \quad
    D_k(\mathbf{b}+\mathbf{c}) \in \mathbf{C}^{(k-1)}.
    \]
    By lemma~\ref{lemma:bpe-decompression-append},
    \[
    D_k(\mathbf{a}+\mathbf{b}) = D_k(\mathbf{a}) + D_k(\mathbf{b}),
    \qquad
    D_k(\mathbf{b}+\mathbf{c}) = D_k(\mathbf{b}) + D_k(\mathbf{c}).
    \]
    Applying the inductive hypothesis at stage $k-1$ yields
    \[
    D_k(\mathbf{a}) + D_k(\mathbf{b}) + D_k(\mathbf{c}) \in \mathbf{C}^{(k-1)}.
    \]
    Using lemma~\ref{lemma:bpe-decompression-append} once more,
    \[
    D_k(\mathbf{a}+\mathbf{b}+\mathbf{c}) \in \mathbf{C}^{(k-1)}.
    \]

    It remains to prove that $\mathbf{a}+\mathbf{b}+\mathbf{c}$ is stage-$k$ reduced.
    Because $\mathbf{b} \neq \boldsymbol{\epsilon}$, the boundary between $\mathbf{a}$ and $\mathbf{c}$
    does not create an adjacent pair in $\mathbf{a}+\mathbf{b}+\mathbf{c}$.
    Every adjacent pair of tokens in $\mathbf{a}+\mathbf{b}+\mathbf{c}$ lies either in the prefix
    $\mathbf{a}+\mathbf{b}$ or in the suffix $\mathbf{b}+\mathbf{c}$.
    Since both of those sequences are stage-$k$ reduced, so is $\mathbf{a}+\mathbf{b}+\mathbf{c}$.

    A final application of theorem~\ref{thm:bpe-canonical-characterization} gives
    \[
    \mathbf{a}+\mathbf{b}+\mathbf{c} \in \mathbf{C}^{(k)},
    \]
    as required.
\end{proof}

\begin{corollary}[BPE last-pair extension]
    \label{cor:bpe-last-pair-extension}
    Assume the tokenization function of this chapter is induced by a BPE merge system, so that
    $\mathbf{C} = \mathbf{C}^{(N)}$.
    Let $\mathbf{a} \in \mathbf{C}$ be nonempty and let $\boldsymbol{\tau}$ be a token.
    If
    \[
    \psi(\mathbf{a}[-1], \boldsymbol{\tau}),
    \]
    then
    \[
    \mathbf{a}+\boldsymbol{\tau} \in \mathbf{C}.
    \]
\end{corollary}
\begin{proof}
    Let
    \[
    \mathbf{a}' := \mathbf{a}[:-1],
    \qquad
    \mathbf{b} := [\mathbf{a}[-1]],
    \qquad
    \mathbf{c} := [\boldsymbol{\tau}].
    \]
    Then
    \[
    \mathbf{a}'+\mathbf{b} = \mathbf{a} \in \mathbf{C}.
    \]
    Also, in the BPE specialization, the condition
    \[
    \psi(\mathbf{a}[-1], \boldsymbol{\tau})
    \]
    is exactly the statement that
    \[
    \mathbf{b}+\mathbf{c} \in \mathbf{C}.
    \]
    Applying theorem~\ref{thm:bpe-overlap-concatenation} with $k=N$ gives
    \[
    \mathbf{a}'+\mathbf{b}+\mathbf{c} \in \mathbf{C},
    \]
    i.e.
    \[
    \mathbf{a}+\boldsymbol{\tau} \in \mathbf{C}.
    \]
\end{proof}
